\documentclass[12pt,a4paper]{ctexart}
\usepackage[margin=1in]{geometry}
\usepackage{graphicx}
\usepackage[dvipsnames]{xcolor}
\usepackage{amsmath}
\usepackage{amssymb}
\usepackage{booktabs}
\usepackage{array}
\usepackage{enumitem}
\usepackage{listings}
\usepackage{hyperref}

\definecolor{codebg}{HTML}{F7F8FA}
\definecolor{codeframe}{HTML}{D9DEE7}
\definecolor{codekeyword}{HTML}{0066CC}
\definecolor{codestring}{HTML}{A31515}
\definecolor{codecomment}{HTML}{008000}

\lstdefinestyle{reportcode}{
  backgroundcolor=\color{codebg},
  basicstyle=\ttfamily\small,
  frame=single,
  rulecolor=\color{codeframe},
  frameround=ffff,
  columns=fullflexible,
  breaklines=true,
  breakatwhitespace=true,
  tabsize=4,
  showstringspaces=false,
  keepspaces=true,
  keywordstyle=\color{codekeyword}\bfseries,
  commentstyle=\color{codecomment},
  stringstyle=\color{codestring},
  numbers=left,
  numberstyle=\tiny\color{gray},
  numbersep=8pt,
  xleftmargin=1.8em,
  framexleftmargin=1.2em,
  aboveskip=0.8em,
  belowskip=0.8em
}

\lstset{style=reportcode}
\hypersetup{
  colorlinks=true,
  linkcolor=MidnightBlue,
  urlcolor=RoyalBlue,
  citecolor=MidnightBlue
}

\title{并行程序设计与算法实验\\实验一：基于 MPI 的并行矩阵乘法}
\author{姓名：元朗曦 \quad 学号：23336294}
\date{\today}

\begin{document}
\maketitle

\section{实验目的}
本实验的目标是使用 MPI 实现通用矩阵乘法的并行计算，并分析不同进程数量和不同矩阵规模下的性能变化。通过本实验希望达到以下目的：
\begin{enumerate}[itemsep=0.3em]
  \item 熟悉 MPI 的基本编程模型，以及主从式并行计算的组织方式。
  \item 掌握矩阵乘法的分块分发、结果汇总与通信开销分析方法。
  \item 理解矩阵规模、进程数量与程序运行时间之间的关系。
  \item 思考在内存受限和稀疏矩阵场景下，如何进一步设计更高效的矩阵乘法算法。
\end{enumerate}

\section{实验原理}
设矩阵 $A \in \mathbb{R}^{m \times n}$，$B \in \mathbb{R}^{n \times k}$，则乘积矩阵 $C \in \mathbb{R}^{m \times k}$ 的定义为
\begin{equation}
C_{i,j} = \sum_{p=1}^{n} A_{i,p} B_{p,j}.
\end{equation}

经典三重循环矩阵乘法的时间复杂度为 $O(mnk)$，对应的浮点运算量可近似为
\begin{equation}
\text{FLOPs} \approx 2mnk.
\end{equation}

若程序运行时间为 $t$ 秒，则浮点性能可表示为
\begin{equation}
\text{GFLOPS} = \frac{2mnk}{t \times 10^9}.
\end{equation}

本实验采用 MPI 对矩阵乘法进行并行化。其基本思路如下：
\begin{enumerate}[itemsep=0.3em]
  \item 由 0 号进程随机生成矩阵 $A$ 和 $B$。
  \item 将矩阵 $A$ 按行划分为若干连续块，分发到各个进程。
  \item 将矩阵 $B$ 广播到所有进程，使每个进程都能访问完整的右矩阵。
  \item 各进程独立计算自己负责的 $A$ 子块与 $B$ 的乘积，得到局部结果块。
  \item 将局部结果块回收至 0 号进程，拼接成最终矩阵 $C$。
\end{enumerate}

这种实现的优点是结构清晰、易于实现，适合验证 MPI 并行矩阵乘法的基本性能特征。其主要开销来自两部分：一是矩阵块分发与回收的通信开销，二是各进程本地乘法计算的时间。随着进程数量增加，计算部分通常下降，但通信与同步开销会逐渐上升，因此加速比不会无限线性增长。

\section{实验内容}
本实验在 \texttt{lab/01/src} 中完成，核心程序为 MPI 矩阵乘法实现，编译后生成可执行文件 \texttt{matmul\_mpi}。程序输入为 $m,n,k$ 三个整数，矩阵元素随机生成，默认输出矩阵 $A$、$B$、$C$ 以及计算时间和 GFLOPS。

实验内容包括以下几部分：
\begin{enumerate}[itemsep=0.3em]
  \item 实现 MPI 版本的通用矩阵乘法程序。
  \item 使用 \texttt{mpirun} 在不同进程数下运行程序，观察运行时间变化。
  \item 在不同矩阵规模下测试程序性能，分析规模扩大对耗时和吞吐率的影响。
  \item 回答实验要求中的讨论题，并总结当前实现的局限性。
\end{enumerate}

实验运行方式如下：
\begin{lstlisting}[language=bash]
cd lab/01/src
make
mpirun --oversubscribe -np 4 ./matmul_mpi 512 512 512 --no-print
python3 benchmark.py \
    --sizes 128,256,512,1024,2048 \
    --processes 1,2,4,8,16 \
    --skip-build \
\end{lstlisting}

\subsection*{实验结果}

\begin{table}[htbp]
  \centering
  \renewcommand{\arraystretch}{1.25}
  \resizebox{\textwidth}{!}{%
  \begin{tabular}{c|c|c|c|c|c}
    \toprule
    Size & np=1 & np=2 & np=4 & np=8 & np=16 \\
    \midrule
    128 & 0.003445 s & 0.004889 s & 0.003381 s & 0.002468 s & 0.002539 s \\
    256 & 0.004088 s & 0.003379 s & 0.002596 s & 0.009158 s & 0.003148 s \\
    512 & 0.023330 s & 0.013122 s & 0.008196 s & 0.009276 s & 0.012584 s \\
    1024 & 0.206837 s & 0.127166 s & 0.114528 s & 0.120935 s & 0.203254 s \\
    2048 & 3.262171 s & 2.246921 s & 1.327422 s & 1.273556 s & 1.131457 s \\
    \bottomrule
  \end{tabular}%
  }
\end{table}

表格如下。由于矩阵乘法属于计算密集型任务，当矩阵规模较小时，MPI 通信开销和进程启动开销会比较显著；当规模增大后，计算量增长更快，程序更容易体现并行带来的收益。与此同时，如果进程数继续增加而问题规模不变，每个进程分到的工作会变少，此时通信与同步占比会上升，导致效率下降。

\subsection*{讨论题回答}
\textbf{1. 在内存有限的情况下，如何进行大规模矩阵乘法计算？}

当矩阵无法一次性完整装入内存时，应采用分块计算、外存计算或流式计算策略。基本思想是把大矩阵切分成多个较小的子块，每次只加载参与当前计算的块进入内存，计算局部乘积后立即写回结果。这样可以显著降低峰值内存占用。

在并行环境中，常见做法包括：
\begin{enumerate}[itemsep=0.3em]
  \item \textbf{二维分块}：把 $A$、$B$、$C$ 都切成子块，让多个进程协同计算不同块。
  \item \textbf{分阶段广播}：不一次性广播完整矩阵，而是按列块或行块分批传输。
  \item \textbf{外存分块}：当内存不足时，将中间块缓存到磁盘，采用多轮 I/O 与计算交替进行。
  \item \textbf{通信与计算重叠}：在加载下一批数据的同时计算当前批次，减少空等时间。
\end{enumerate}

对于极大规模矩阵，合理的块大小选择非常关键：块太小会导致通信和调度开销过大，块太大又会增加内存占用并降低缓存命中率。实际实现中通常需要结合机器内存容量、缓存层次结构和网络带宽进行权衡。

\textbf{2. 如何提高大规模稀疏矩阵乘法性能？}

稀疏矩阵乘法的关键不在于减少循环层数，而在于尽量避免对零元素进行无效计算。对于稀疏矩阵，若仍使用普通稠密矩阵乘法，会浪费大量时间和内存。因此通常需要采用稀疏存储格式和稀疏算法。

可行的优化方向包括：
\begin{enumerate}[itemsep=0.3em]
  \item \textbf{采用稀疏存储格式}：如 CSR、CSC、COO 等，只存储非零元素及其索引。
  \item \textbf{基于非零结构进行计算}：乘法时只遍历非零项，并利用索引快速定位对应元素。
  \item \textbf{负载均衡}：稀疏矩阵中各行/各块的非零数可能差异很大，MPI 划分时不能只按行数平均分配，而应按非零元数量或预计计算量进行分配。
  \item \textbf{局部性优化}：对非零结构进行重排或分块，以提高缓存命中率和向量化效率。
  \item \textbf{通信压缩}：并行时只传输非零块或必要索引，减少网络通信量。
\end{enumerate}

因此，稀疏矩阵乘法的核心优化目标是“跳过零计算、减少无效数据搬运、保持负载均衡”。如果稀疏模式高度不规则，通常还需要结合图划分、超图划分或动态任务调度进一步优化。

\subsection*{结果分析}
本实验的实现属于基础的行块并行方案，优点是简单稳定，适合验证 MPI 并行矩阵乘法的基本流程；不足之处在于 $B$ 被完整广播到所有进程，存在一定的内存重复占用和通信开销。对于更大规模问题或更高进程数场景，可以进一步采用二维块分布、SUMMA 类算法或分层通信方案，以降低广播和回收成本。

\section{实验总结}
通过本实验，我完成了基于 MPI 的并行通用矩阵乘法实现，并理解了并行矩阵乘法的基本设计方法。实验表明，矩阵乘法的性能不仅取决于计算量本身，还受到数据分布方式、通信开销和负载均衡的显著影响。

本实验的主要收获如下：
\begin{enumerate}[itemsep=0.3em]
  \item MPI 并行程序需要同时考虑通信和计算，不能只关注局部计算效率。
  \item 对于矩阵乘法这类计算密集型任务，矩阵划分方式直接影响并行效率。
  \item 在内存受限和稀疏矩阵场景下，算法设计思路需要从“直接乘法”转向“分块、压缩和结构化计算”。
\end{enumerate}

总体而言，本实验帮助我建立了从串行优化过渡到并行优化的思维方式，也为后续进一步学习更高效的并行矩阵乘法算法打下了基础。

\end{document}
